% HRCP Research Proposal — Fall 2026 term
% Ben Charette · Harvard College · Department of Physics / IceCube
% Build: pdflatex hrcp_proposal_fall2026.tex
\documentclass[12pt,letterpaper]{article}
\usepackage[margin=1in]{geometry}
\usepackage{setspace}
\usepackage[T1]{fontenc}
\usepackage{times}
\usepackage{microtype}
\usepackage{titlesec}
\usepackage{enumitem}
\usepackage[hidelinks]{hyperref}

\titlespacing*{\section}{0pt}{8pt}{2pt}
\titleformat{\section}{\normalfont\bfseries}{\thesection.}{0.6em}{}
\setlength{\parskip}{0pt}
\setlength{\parindent}{0.4in}

\begin{document}

\begin{center}
{\large\bfseries Muon-Induced Background in the Unanalyzed Years}\\[2pt]
{\large\bfseries of the DM-Ice17 Dark Matter Detector}\\[8pt]
{\normalsize Ben Charette \quad$\cdot$\quad September 2026}
\end{center}

\vspace{2pt}
\doublespacing

\section{Background and State of the Field}

Dark matter makes up roughly a quarter of the universe's energy content, yet it has
never been detected in a laboratory. For more than two decades one experiment has
claimed otherwise. DAMA/LIBRA, operating beneath the Gran Sasso mountain in Italy, uses
crystals of thallium-doped sodium iodide, NaI(Tl), a material that emits a brief flash
of light when a particle deposits energy in it. DAMA reports that its rate of very
low-energy flashes rises and falls once per year, peaking near June~2, and interprets
this as the Earth's orbital motion carrying its detector through a galactic halo of
dark matter particles: in June the Earth moves into the oncoming wind of dark matter,
in December it moves with it, and the collision rate should follow [1].

No other experiment has confirmed the result, and the central objection is mundane.
Laboratories also have seasons. Temperature, humidity, radon concentration, and the
flux of cosmic-ray muons through the rock all vary annually, and any of them could
imprint a yearly cycle on a detector. Distinguishing a galactic signal from a
terrestrial one therefore comes down to \emph{phase}. A dark matter signal is set by
the Earth's orbit and keeps its June maximum everywhere on the planet, while most
environmental effects follow the local seasons and invert in the southern hemisphere.

DM-Ice17 was built to exploit exactly that asymmetry [2]. It consists of two 8.5\,kg
NaI(Tl) crystals deployed in 2010 at the bottom of two strings of the IceCube Neutrino
Observatory and frozen permanently into the Antarctic ice at 2450\,m depth. IceCube
instruments a cubic kilometre of that ice with thousands of light sensors in order to
detect neutrinos [3]; muons---heavy, short-lived relatives of the electron, produced
when cosmic rays strike the upper atmosphere---penetrate deeply enough to reach the
array constantly and are IceCube's dominant background. For DM-Ice17 this arrangement
is a gift: it is the only NaI dark matter detector ever operated inside an active
cubic-kilometre muon telescope, so the particles passing through its crystals can be
identified independently by the surrounding array.

The published DM-Ice muon analysis established three things [4]. Muons cross each
crystal about 2.93 times per day; that rate modulates by $12.3 \pm 1.7\%$ and peaks
around January~22, in austral summer, four months from the DAMA phase; and muons
leave behind a faint afterglow in the crystal, phosphorescence with a decay time of
$5.5 \pm 0.5$\,s, which the authors characterized as noise and argued cannot reproduce
the DAMA signal. Crucially, that analysis used 3.5 years of data and stops in January
2015, two months after the crystals' photomultiplier gain was raised roughly tenfold.
Twelve years of data were taken. Fewer than four have ever been analyzed, and the
unanalyzed years are the ones the modern dark matter search would use.

\section{Objective, Significance, and Implications}

My objective is to extend the characterization of muons and muon-induced background
into the unanalyzed 2015--2020 high-gain era, and to convert the known existence of
afterglow into a quantitative statement about a dark matter search: \textbf{what
fraction of events in the 2--6\,keV electron-equivalent (keVee) signal region is
muon-induced, what fraction of that an IceCube veto removes, and at what cost in
lost exposure.}

This goes beyond the published work in two specific ways. First, coverage: the
tenfold gain change did not merely shift a calibration, it inverted the behavior of
the pulse-shape parameter used to identify muons, which is why no selection was ever
published for those years. I have already solved that problem, so the era is now
open for the first time. Second, framing: the published measurement established that
phosphorescence exists and decays in about 5.5 seconds, but treated it as a noise
source rather than asking how much of it survives a low-energy event selection and
lands inside the signal region. Those are different questions, and only the second
one bears on whether a June-phase signal could be environmental.

I want to be direct about what this cannot do. With 17\,kg of crystal, DM-Ice17 will
never outmeasure DAMA/LIBRA's roughly 250\,kg at a lower threshold, and this project
will not produce a competitive dark matter limit. Its value is that DM-Ice17 is the
only place where the muon-induced contribution to a NaI signal region, and the
effectiveness of an active veto against it, can be \emph{measured} rather than
simulated---numbers that feed directly into the successor experiment, COSINE-100 [5],
and into any future veto-instrumented NaI deployment.

\section{Research Plan and Methodology}

\noindent\textbf{Preparation already completed.} This is a continuation of work I
began last summer, and its infrastructure is built and validated. I reduced the full
2011--2022 dataset---2.38 billion recorded events---into a compact columnar form, and
built the first true live-time database for DM-Ice in the group's code base, covering
95{,}174 runs. Applying the published method to the low-gain era, I recover a muon rate
of $2.918 \pm 0.038$ per crystal per day and a flux modulation of $10.2 \pm 1.6\%$
peaking on day $30 \pm 9$, all consistent within uncertainties with the published
$2.93 \pm 0.04$, $12.3 \pm 1.7\%$, and January~22 [4]---my evidence that the pipeline
is trustworthy. I then diagnosed why a naive transfer of that selection to the
high-gain era overcounted muons by a factor of 751, and fixed it with a discriminant
not previously used on these data, reaching 0.95 efficiency and 0.93 purity.

\noindent\textbf{Step 1 (weeks 1--2): establish the low-energy threshold.} My own
calibration notes flag that a feature previously identified as a 46.5\,keV lead-210
line is more likely the detector's trigger turn-on, its width being far too narrow for
a real line. Because everything downstream depends on it, I will first measure where
the turn-on actually lies and define the lowest defensible analysis band from the data
rather than from the literature.

\noindent\textbf{Step 2 (weeks 3--6): measure afterglow leakage into the signal
region.} For each of roughly a thousand tagged muons per year per detector, I will
count low-energy events in delayed windows following the muon and compare them against
randomized control windows drawn from the same runs. Because both window classes are
drawn from identical live-time, the comparison is self-normalizing and immune to the
live-time uncertainties that dominate any rate measurement. This yields the
muon-induced fraction of the signal-region rate and its decay profile, which I can
check against the published 5.5\,s constant as a validation.

\noindent\textbf{Step 3 (weeks 7--9): veto efficiency and its cost.} Using the muon
sample tagged by IceCube, I will compute what fraction of that leakage an active veto
removes as a function of the veto window length, together with the exposure sacrificed.
At three muons per crystal per day, even a 30-second window costs roughly a tenth of a
percent of exposure, so I expect the interesting result to be how much background it
buys.

\noindent\textbf{Step 4 (weeks 10--13): extend to the high-gain era.} Low-energy
events in the high-gain years are digitized on a channel my current reduction does not
carry, so extending Steps 2 and 3 past 2015 requires a re-extraction that I have
written and tested on a subsample. This is the novel-coverage half of the project and I
will launch it in week 1 to run unattended, but I have deliberately placed it after a
complete low-gain result so that a storage bottleneck cannot cost me the term.

\noindent\textbf{Step 5 (weeks 14--15): one open systematic, and write-up.} The 2015
muon rate sits well below 2016--2020 and is not a selection artifact; I will test
whether per-run live-time is overestimated because it cannot resolve dead time within a
run. I will then write an internal collaboration note and draft a conference abstract.

\section{Timeline}

The Harvard term runs from early September to mid-December, and I can commit 10 to 12
hours per week. Long extractions run unattended on IceCube's Cobalt and NPX computing
systems, so the heaviest computation does not compete with coursework. The five steps
above are scheduled across weeks 1--2, 3--6, 7--9, 10--13, and 14--15 respectively,
with two review points: a measured threshold at the end of week 2, and the
muon-induced leakage fraction at the end of week 6, each discussed with my mentor
before I build on it. I will present Steps 2 and 3 to the group in early November. The
project is deliberately ordered so that a complete, self-contained result exists by
week 9 even if the high-gain extension in Step 4 does not finish.

\section{Faculty Involvement and Support Network}

My faculty mentor is Dr.\ Carlos Arg\"uelles-Delgado, who leads Harvard's IceCube group
and has supervised this project throughout. I attend his weekly group meetings on
Tuesdays from 3--5pm, where the state of the whole lab is discussed, and I will meet
with him one-on-one at least four times over the term, including at the two review
points above. He sets the physics priorities, and he has previously required me to
freeze and defend a completed analysis in writing before moving on---a practice that
produced the memo this proposal is built from. He also provides the access the project
depends on: IceCube's Cobalt and NPX systems, the group's allocation on the FAS
research cluster, and the IceCube Collaboration channel through which any result of
mine receives internal review before publication.

For day-to-day technical guidance I have weekly check-ins, at minimum, with Dr.\ Felix
Yu, a postdoc in the group whose work on reconstruction methods overlaps my own and who
defined several of the tasks I completed last summer; I will adjust that frequency to my
needs as the term progresses. Working in the department's Laboratory for Particle
Physics and Cosmology also puts me among the other students on the team, which has been
the most reliable way I have found to learn techniques that are not written down
anywhere. One question I will bring to my mentor early is whether to frame the term's
output primarily as a background-characterization result for the collaboration or as a
modulation measurement, since that affects how much weight Step 4 should carry.

\section{Citations}
\begin{spacing}{1.0}
\begin{enumerate}[leftmargin=*,itemsep=3pt]
\item Bernabei, R., et al. (DAMA/LIBRA Collaboration). (2018). First model independent
results from DAMA/LIBRA-phase2. \emph{Universe}, 4(11), 116.
doi: 10.3390/universe4110116

\item Cherwinka, J., et al. (DM-Ice Collaboration). (2014). First data from DM-Ice17.
\emph{Physical Review D}, 90, 092005. doi: 10.1103/PhysRevD.90.092005

\item Aartsen, M. G., et al. (IceCube Collaboration). (2017). The IceCube Neutrino
Observatory: instrumentation and online systems. \emph{Journal of Instrumentation},
12, P03012. doi: 10.1088/1748-0221/12/03/P03012

\item Cherwinka, J., et al. (DM-Ice Collaboration). (2016). Measurement of muon annual
modulation and muon-induced phosphorescence in NaI(Tl) crystals with DM-Ice17.
\emph{Physical Review D}. arXiv:1509.02486

\item Adhikari, G., et al. (COSINE-100 Collaboration). (2018). An experiment to search
for dark-matter interactions using sodium iodide detectors. \emph{Nature}, 564, 83--86.
doi: 10.1038/s41586-018-0739-1

\item Hubbard, A. J. F. (2015). \emph{Muon-Induced Backgrounds in the DM-Ice17 NaI(Tl)
Dark Matter Detector}. PhD thesis, University of Wisconsin--Madison.

\item Tilav, S., et al. (IceCube Collaboration). (2010). Atmospheric variations as
observed by IceCube. arXiv:1001.0776
\end{enumerate}
\end{spacing}

\end{document}
